\section{信源分类与单符号信源}
\warmth{0}\quad\whisper{先看一次输出一个随机符号的信源，这是后面序列信源的基本积木。}

\begin{strand}
信源可以按输出的对象来分：如果每次输出一个符号，叫单符号信源；
如果连续输出一串符号，叫消息序列信源。也可以按取值集合分为离散信源和连续信源。
\end{strand}

\block{分类地图}
\noindent\begin{tabularx}{\linewidth}{@{}l L@{}}
\toprule
{\color{indigo}分类角度} & \multicolumn{1}{l}{\color{indigo}常见类型}\\
\midrule
按取值集合 & 离散信源，也可称数字信源；连续信源，也可称模拟信源。\\
按输出结构 & 单符号信源；消息序列信源，也就是随机过程形式的信源。\\
\bottomrule
\end{tabularx}

\begin{definition}[离散单符号信源]
设离散单符号信源为随机变量 \(X\)，其取值集合为
\[
\Xcal=\{x_1,x_2,\cdots,x_n\}.
\]
若每个符号 \(x_i\) 出现的概率为 \(\pmf(x_i)\)，则该信源的统计描述为
\[
\left(
\begin{array}{cccc}
x_1 & x_2 & \cdots & x_n\\
\pmf(x_1) & \pmf(x_2) & \cdots & \pmf(x_n)
\end{array}
\right),
\]
其中
\[
0\leq \pmf(x_i)\leq 1,\qquad
\sum_{i=1}^{n}\pmf(x_i)=\answerin[1.2cm]{1}.
\]
\end{definition}

\begin{definition}[连续单符号信源]
若信源输出 \(X\) 是连续随机变量，例如 \(X\in(a,b)\)，则通常用概率密度函数
\(\pmf(x)\) 描述。这里的 \(\pmf(x)\) 不是某一点的概率；当 \(dx\) 很小时，
\[
\pmf(x)\,dx \approx P\{x\leq X<x+dx\}.
\]
更一般地，
\[
P\{a\leq X\leq b\}=\int_a^b \pmf(x)\,dx.
\]
也就是说，\(\pmf(x)\) 描述的是单位长度小区间中的
\answerin[3.4cm]{概率密度}。
\end{definition}

\begin{remark}
离散情形用“概率质量”看每个符号，连续情形用“概率密度”看区间附近。
二者的共同点是：都在回答信源输出落在某些取值上的可能性。
本章聚焦离散信源；连续信源及其微分熵将在后续课程中展开。
\end{remark}

\begin{yourturn}
判断下面两个对象分别属于哪类信源：
\[
\text{二进制符号 } \{0,1\}:\ \answerin[3cm]{离散信源},
\qquad
\text{模拟电压 } X(t):\ \answerin[3cm]{连续信源}.
\]
\workspace[2]
\end{yourturn}

\recall{为什么连续信源中 \(\pmf(x)\) 通常叫概率密度而不是概率？}
